\titleframe

% Contoh Kode Slide
%\section{Introduction}
%\begin{frame}{\insertsectionhead}
%	\framesubtitle{A short introduction to Trigon}
%	\themename is a modern, elegant and versatile theme for Beamer, inspired by
%	the
%	\href{https://github.com/matze/mtheme}{\textsc{metropolis} theme} from Matthias
%	Vogelgesang.
%	\vfill
%	\themename comes with lots of nice extra features
%	\begin{itemize}
%		\item Multiple style variations for title, section and normal slides
%		\item Simple customization of theme colors
%		\item Lots of convenient options to tweak the design
%	\end{itemize}
%\end{frame}

% Contoh Kode Tabel
%  \begin{table}[H]
% \centering
% \caption{A nice table example}
% \begin{tabular}{@{} lccc @{}}
%    \toprule
%    & \textbf{Velocity} & \textbf{Angle}  & \textbf{Vertical force} \\
%    & $U$ & $\alpha$  & $F_z$ \\
%    & [m/s] & [$^\circ$]  & [N] \\
%    \midrule
%    2D simulation  & 9 & 2 & 9.23 \\
%    3D simulation  & 10.0 & 3 & 15.039 \\
%    Experiment A   & 11.31 & 2.5 & 13.2 \\
%    Experiment B   & 11.26 & 2.7 & 12.6 \\
%    Experiment C   & 11.33 & 2.47 & 13.6 \\
%    \bottomrule
%  \end{tabular}
%\end{table}

% Contoh Kode Gambar
%  \begin{figure}[ht!]
%	\begin{subfigure}[b]{0.3\textwidth}
%		\frame{\includegraphics[width=\textwidth]{layout_example-03.jpg}}
%		\caption*{plain}
%	\end{subfigure}
%	\hspace{\fill}
%	\begin{subfigure}[b]{0.3\textwidth}
%		\frame{\includegraphics[width=\textwidth]{layout_example-02.jpg}}
%		\caption*{style1}
%	\end{subfigure}
%	\hspace{\fill}
%	\begin{subfigure}[b]{0.3\textwidth}
%		\frame{\includegraphics[width=\textwidth]{layout_example-01.jpg}}
%		\caption*{style2 (default)}
%	\end{subfigure}
%\end{figure}

% Contoh Kode Blocks
% \begin{block}{Regular block}
%	Just a regular block
%\end{block}
%\begin{alertblock}{Alert block}
%	Some important thing
%\end{alertblock}
%\begin{exampleblock}{Example block}
%	No difference with regular block to avoid excessive distraction
%\end{exampleblock}

\setbeamertemplate{frame footer}{Alauddin Maulana Hirzan, M. Kom - Universitas Semarang}

\section{\textit{Internet of Things} dan Aplikasinya}
\begin{frame}{\insertsectionhead}
	\framesubtitle{\textit{Internet of Things} dan Aplikasi \#1}
	\justifying
	Tujuan IoT tidak hanya menghubungkan hal-hal seperti mesin, perangkat dan peralatan, tetapi juga memungkinkan hal-hal untuk berkomunikasi, bertukar data kontrol dan informasi lain yang diperlukan saat menjalankan aplikasi menuju tujuan mesin.
	\vfill
	Perangkat ini dapat berupa perangkat pengumpul data tempat berbagai sensor terpasang seperti suhu, kelembaban, cahaya, dll. Atau dapat berupa perangkat penggerak data tempat aktuator terhubung seperti relai.
	\vfill
	Sehingga dapat diimplementasikan sesuai dengan kebutuhan yang ada
\end{frame}

\begin{frame}{\insertsectionhead}
	\framesubtitle{\textit{Internet of Things} dan Aplikasi \#2}
	\justifying
	Akan ada 26 miliar perangkat yang terhubung ke internet of things pada tahun 2020 menurut Gartner Inc. Aplikasi IOT tidak terbatas pada beberapa kelompok area, tetapi dapat diterapkan ke berbagai domain baik individu, industri, atau infrastruktur.
	\vfill
	Produk IoT diklasifikasikan berdasarkan jenis aplikasinya. Ini bisa berupa \textit{smart wearable, smart city, smart home, smart enterprise, atau smart environment}.
\end{frame}

\begin{frame}{\insertsectionhead}
	\framesubtitle{\textit{Internet of Things} dan Aplikasi \#2}
	\justifying
	Beberapa area penerapan IoT dibahas di bawah ini:
	\vfill
	\begin{itemize}
		\item Otomatisasi Rumah dan Gedung
		\item Industri
		\item Energi
		\item Sistem Medis dan Perawatan Kesehatan
		\item Angkutan / Transportasi
		\item Lingkungan
	\end{itemize}
\end{frame}

\begin{frame}{\insertsectionhead}
	\framesubtitle{\textit{Internet of Things} dan Aplikasi \#3}
	\justifying
	 \begin{figure}[ht!]
			\begin{subfigure}[b]{0.5\textwidth}
					\frame{\includegraphics[width=\textwidth]{market-share.jpg}}
				\end{subfigure}
		\end{figure}
\end{frame}

\begin{frame}{\insertsectionhead}
	\framesubtitle{\textit{Internet of Things} dan Aplikasi \#4}
	\justifying
	\begin{figure}[ht!]
		\begin{subfigure}[b]{0.7\textwidth}
			\frame{\includegraphics[width=\textwidth]{indo-market.jpg}}
		\end{subfigure}
	\end{figure}
\end{frame}

\begin{frame}{\insertsectionhead}
	\framesubtitle{\textit{Internet of Things} dan Aplikasi \#5}
	\justifying
	Dikutip dari \textbf{nextren.grid.id}, Gunawan mengatakan bahwa:
	\vfill
	\begin{quote}
		\justifying
		Secara mendetil dijelaskan kalau kontribusi peningkatan itu nantinya akan terjadi di industri aplikasi sebesar 43 persen. Lalu selanjutnya akan ada dari industri platform (35 persen), devices (13 persen), dan network (9 persen). Terkait alasannya, Gunawan pun menyebut kalau pertumbuhan aplikasi itu didasari oleh awalan dari pertumbuhan devices yang akan diikuti setelahnya. "Akan banyak di industri aplikasi terlebih dahulu karena untuk di device akan banyak permintaan teknologi aliansi bisa membangun perangkat disini," tuturnya.
	\end{quote}
\end{frame}

\begin{frame}{\insertsectionhead}
	\framesubtitle{Aplikasi IoT -  Otomatisasi Rumah dan Gedung \#1}
	\justifying
	Perangkat IOT digunakan untuk memantau dan mengontrol sistem elektronik, listrik dan mekanik di rumah dan gedung untuk meningkatkan kenyamanan dan keamanan. Tugas IOT dalam domain ini meliputi:
	\vfill
	\begin{itemize}
		\item Pencahayaan cerdas dengan mengadaptasi peralihan berdasarkan kondisi sekitar
		\item Aplikasi web dan aplikasi seluler mengaktifkan lampu terhubung nirkabel dan internet
		\item Manajemen dan kontrol peralatan pintar
		\item Sistem deteksi intrusi, sistem alarm, dan sistem pengawasan
		\item Sistem keamanan seperti deteksi asap dan gas
		\item Manajemen hiburan rumah seperti video, audio, dan proyektor
	\end{itemize}
\end{frame}

\begin{frame}{\insertsectionhead}
	\framesubtitle{Aplikasi IoT -  Otomatisasi Rumah dan Gedung \#2}
	\justifying
	\begin{figure}[ht!]
		\begin{subfigure}[b]{0.7\textwidth}
			\frame{\includegraphics[width=\textwidth]{smart-home.jpg}}
		\end{subfigure}
	\end{figure}
\end{frame}

\begin{frame}{\insertsectionhead}
	\framesubtitle{Aplikasi IoT -  Otomatisasi Rumah dan Gedung \#3}
	\justifying
	Banyak kemungkinan yang bisa dilakukan untuk perangkat rumah pintar, dan otomatisasi rumah menjadi impian di masa depan. Para peneliti telah memperkirakan bahwa akan ada pertumbuhan luar biasa dalam pengiriman peralatan rumah pintar. 
	\vfill
	Hal ini termasuk semua peralatan pintar (mesin cuci, pengering, lemari es, dll.), bersama dengan perangkat keselamatan dan keamanan rumah pintar dan sistem keamanan (sensor, monitor, kamera, sistem alarm pencuri, dan peralatan hemat energi rumah pintar, seperti termostat dan pencahayaan pintar) .
\end{frame}

\begin{frame}{\insertsectionhead}
	\framesubtitle{Aplikasi IoT -  Industri \#1}
	\justifying
	IOT berkaitan dengan optimasi real-time jaringan produksi dan rantai pasokan dalam industri manufaktur dengan sensor jaringan, aktuator, sistem kontrol dan mesin.
	\vfill	
	Sehingga, tugas Industrial Internet of Things (IIOT) adalah:
	\begin{itemize}
		\item Pemantauan dan pengendalian proses secara real-time
		\item Menyebarkan mesin pintar, sensor pintar, dan pengontrol pintar dengan teknologi komunikasi dan internet yang dipatenkan
		\item Maksimalkan keselamatan, keamanan, dan keandalan melalui otomatisasi dan kontrol presisi tinggi
	\end{itemize}
\end{frame}

\begin{frame}{\insertsectionhead}
	\framesubtitle{Aplikasi IoT -  Industri \#2}
	\justifying
	\begin{figure}[ht!]
		\begin{subfigure}[b]{0.7\textwidth}
			\frame{\includegraphics[width=\textwidth]{industry.png}}
		\end{subfigure}
	\end{figure}
\end{frame}

\begin{frame}{\insertsectionhead}
	\framesubtitle{Aplikasi IoT -  Industri \#3}
	\justifying
	"Industri Cerdas" adalah (r)evolusi yang dibangun di atas banyak kemajuan teknologi yang berbeda selama beberapa dekade terakhir (evolusi) tetapi juga mengubah cara fungsi pabrik dan tempat kerja (revolusi). 
	\vfill
	Definisi "Cerdas" berarti melakukan sesuatu dengan lebih efisien, lebih fleksibel, dan dengan cara yang lebih ramah lingkungan. Sehingga manufaktur akan lebih aman bagi orang-orang yang bekerja.
	\vfill
	\begin{block}{Info}
		Di masa depan, banyak industri yang akan bergantung kepada teknologi \textit{Internet of Things}. Sehingga topik ini masih akan diminat hingga tahun-tahun kedepan
	\end{block}
\end{frame}

\begin{frame}{\insertsectionhead}
	\framesubtitle{Aplikasi IoT -  Energi \#1}
	\justifying
	Sektor ini mengharuskan IOT berurusan dengan sistem manajemen energi dalam bentuk \textit{Smart Grid}. IOT memperluas manfaat \textit{Smart Grid} di luar otomatisasi, distribusi, dan pemantauan yang dilakukan oleh utilitas. 
	\vfill
	Sehingga, tugas IoT di bidang energi listrik meliputi:
	\begin{itemize}
		\item Infrastruktur Pengukuran Tingkat Lanjut (AMI) - Meteran Pintar
		\item SCADA (Kontrol Pengawasan dan Akuisisi Data) - Pengendalian, dan Pengawasan Pintar
		\item Inverter Cerdas
		\item Pengoperasian remote control dari perangkat yang mengkonsumsi energi
	\end{itemize}
\end{frame}

\begin{frame}{\insertsectionhead}
	\framesubtitle{Aplikasi IoT -  Energi \#2}
	\justifying
	\begin{figure}[ht!]
		\begin{subfigure}[b]{0.58\textwidth}
			\frame{\includegraphics[width=\textwidth]{smart-grid.jpg}}
		\end{subfigure}
	\end{figure}
\end{frame}

\begin{frame}{\insertsectionhead}
	\framesubtitle{Aplikasi IoT -  Energi \#3}
	\justifying
	Smart grid adalah jaringan nasional terencana yang menggunakan teknologi informasi untuk menyalurkan listrik secara efisien, andal, dan aman. Hal ini bisa disebut dengan "listrik dengan otak", "internet energi", dan "elektronet". 
	\vfill
	Definisi yang lebih komprehensif dari NIST adalah "jaringan modern yang memungkinkan aliran energi dua arah dan menggunakan komunikasi dua arah dan kemampuan kontrol yang akan mengarah pada serangkaian fungsi dan aplikasi baru."
\end{frame}

\begin{frame}{\insertsectionhead}
	\framesubtitle{Aplikasi IoT -  Energi \#4}
	\justifying
	\begin{figure}[ht!]
		\begin{subfigure}[b]{1.0\textwidth}
			\frame{
\includegraphics[width=\textwidth]{Iconet.png}}
			\caption*{Dejavu?}
		\end{subfigure}
	\end{figure}
\end{frame}

\begin{frame}{\insertsectionhead}
	\framesubtitle{Aplikasi IoT -  Sistem Medis dan Perawatan Kesehatan \#1}
	\justifying
	Sistem Medis dan Perawatan Kesehatan berbasis Teknologi Internet of Things telah mempengaruhi tingkat perkembangan tertentu dalam sistem medis digital seluler. Hal ini membantu para dokter mendapatkan beberapa parameter fisiologis tubuh manusia ke pusat layanan medis melalui jaringan sensor multi-parameter yang dapat dipakai manusia. 
	\vfill
	Sehingga, tugas IoT di bidang ini antara lain:
	\begin{itemize}
		\item Pemantauan kesehatan jarak jauh
		\item Sistem pemberitahuan darurat
		\item Sisa obat
		\item Pemantauan bayi
	\end{itemize}
\end{frame}

\begin{frame}{\insertsectionhead}
	\framesubtitle{Aplikasi IoT -  Sistem Medis dan Perawatan Kesehatan \#2}
	\justifying
	\begin{figure}[ht!]
		\begin{subfigure}[b]{0.68\textwidth}
			\frame{\includegraphics[width=\textwidth]{smart-medical.jpg}}
		\end{subfigure}
	\end{figure}
\end{frame}

\begin{frame}{\insertsectionhead}
	\framesubtitle{Aplikasi IoT -  Sistem Medis dan Perawatan Kesehatan \#3}
	\justifying
	Dokter dan pasien dapat berbagi informasi secara real time dari satu layar komputer ke layar komputer lainnya. Dan mereka bahkan dapat melihat dan menangkap bacaan dari perangkat medis di lokasi yang jauh. 
	\vfill
	Dengan menggunakan perangkat lunak telemedicine, pasien dapat menemui dokter untuk diagnosis dan pengobatan tanpa harus menunggu janji. Pasien dapat berkonsultasi dengan dokter di kenyamanan rumah mereka.
\end{frame}

\begin{frame}{\insertsectionhead}
	\framesubtitle{Aplikasi IoT -  Transportasi \#1}
	\justifying
	IoT dapat mengintegrasikan kontrol, pemrosesan informasi, dan komunikasi di berbagai sistem transportasi dengan membangun interkoneksi antara pengemudi atau pengguna dan kendaraan. 
	\vfill
	Tugas IoT di bidang transportasi antara lain:
	\vfill
	\begin{multicols}{2}
		\begin{itemize}
			\item Kontrol lalu lintas yang cerdas
			\item Navigasi otonom tak berawak
			\item Komunikasi antar dan intra kendaraan
			\item Respons Darurat Otomatis
			\item Sistem pembayaran tol elektronik
			\item Bantuan keselamatan dan jalan
			\item Parkir cerdas
		\end{itemize}
	\end{multicols}
\end{frame}

\begin{frame}{\insertsectionhead}
	\framesubtitle{Aplikasi IoT -  Transportasi \#2}
	\justifying
	\begin{figure}[ht!]
		\begin{subfigure}[b]{0.68\textwidth}
			\frame{\includegraphics[width=\textwidth]{smart-transport.jpg}}
		\end{subfigure}
	\end{figure}
\end{frame}

\begin{frame}{\insertsectionhead}
	\framesubtitle{Aplikasi IoT -  Transportasi \#3}
	\justifying
	Menurut \textbf{Departemen Transportasi AS}:
	\vfill
	\begin{quote}
		“Sistem Transportasi Cerdas (ITS) menerapkan berbagai teknologi untuk memantau, mengevaluasi, dan mengelola sistem transportasi untuk meningkatkan efisiensi dan keselamatan.”
	\end{quote}
\end{frame}

\begin{frame}{\insertsectionhead}
	\framesubtitle{Aplikasi IoT -  Lingkungan \#1}
	\justifying
	IOT di bidang pemantauan lingkungan terdiri dari sensor parameter lingkungan yang berbeda seperti udara, air, tanah, dll dengan aplikasi berbasis cloud sehingga pemantauan dapat dicapai secara efektif. 
	\vfill
	Tugas IOT di bidang ini meliputi:
	\vfill
	\begin{itemize}
		\item Pemantauan cuaca berbasis cloud
		\item Pemantauan kebisingan dan polusi udara
		\item Sistem deteksi kebakaran
		\item Sistem peringatan dini gempa dan tsunami
		\item Pemantauan kondisi tanah
	\end{itemize}
\end{frame}

\begin{frame}{\insertsectionhead}
	\framesubtitle{Aplikasi IoT -  Lingkungan \#2}
	\justifying
	\begin{figure}[ht!]
		\begin{subfigure}[b]{0.9\textwidth}
			\frame{\includegraphics[width=\textwidth]{smart-env.jpg}}
		\end{subfigure}
	\end{figure}
\end{frame}

\setbeamertemplate{background} 
{
	\includegraphics[width=\paperwidth,height=\paperheight]{thank-you.jpg}
}
\begin{frame}[plain]
\end{frame}